\documentclass[11pt]{article}

\usepackage[utf8]{inputenc}
\usepackage[T1]{fontenc}
\usepackage{lmodern}
\usepackage[margin=1in]{geometry}
\usepackage{amsmath,amssymb}
\usepackage{booktabs}
\usepackage{tabularx}
\usepackage{array}
\usepackage{float}
\usepackage[numbers,sort&compress]{natbib}
\usepackage[hidelinks]{hyperref}

\newcommand{\lean}[1]{\texttt{#1}}
\newcommand{\code}[1]{\texttt{#1}}

\title{SHH Revisited: Don't Just Run Models, Reason About Them}
\author{Jeff Shrager}
\date{Draft, November 2026}

\begin{document}
\sloppy
\maketitle

\begin{center}
\fbox{\parbox{0.9\linewidth}{\centering\large\textbf{Draft for Comment}\footnotemark}}
\end{center}
\footnotetext{This is a working draft. I welcome corrections,
  constructive challenges, and especially proposals for
  collaboration. Email me: \texttt{jshrager@gmail.com}.}
\medskip

\begin{abstract}
In 1988, Tad Hogg, Bernardo Huberman and I proposed that the power law of
practice need not be built into a learning mechanism: it could emerge from
the dynamics of search in a problem space that learning makes more densely
connected, and the same dynamics could make practice hurt, a
problem-solving fan effect. The paper supported this with simulations of a
simple ``bit game'' and with an analytic theory based on a random walk.
Here I revisit that work by formalizing it in the Lean proof assistant, so
that what the theory entails is checked by machine rather than read off
curves. Most of the 1988 mathematics survives: the expected-steps formula
is correct exactly as printed, and its drift and diffusion regimes hold,
one of them in a sharper form. But two of the paper's conclusions do not
follow from its own theory. With no decision procedure, the idealized
model shows a fan effect that the paper says is absent; and with a perfect
decision procedure the ``power-law decay'' it derives is logarithmic, a
curve with every qualitative signature of a power law over the range
plotted, and provably not one over long ranges. In the bit game itself,
exact results reproduce the 1988 simulations' pattern, which locates the
first discrepancy precisely: in the idealization that connects the game to
the theory, not in the game. None of these findings needed new
mathematics. They needed the theory written down precisely enough that its
consequences could be checked.
\end{abstract}

\section{Introduction}
\label{sec:intro}

The power law of practice is one of the most robust regularities in the
psychology of learning: the time to perform a task falls as a power of the
number of times it has been performed \citep{snoddy1926,crossman1959,newell1981}.
In 1988 the leading accounts obtained it in quite different ways. ACT*
obtained it from a strengthening mechanism that itself follows a power law
\citep{anderson1983}; Newell and Rosenbloom obtained it from the chunking
of subgoals \citep{newell1981,rosenbloom1986}. Tad Hogg, Bernardo Huberman
and I proposed a third account \citep{shh1988}. Problem solving is search
through a problem-space graph; learning adds edges to the graph (new
operators) and can improve the procedure that chooses among them. Neither
change mentions a power law. Yet their interaction, under particular
conditions, yields curves that look like one, and under other conditions
yields the opposite: learning that makes performance worse. We called that
the problem-solving fan effect, by analogy with the memory fan effect
\citep{anderson1974}, and connected it to the utility problem in machine
learning \citep{minton1985,markovitch1988}.

That paper did its reasoning in two ways. It simulated a small task, the
bit game, and it derived an analytic theory, a one-dimensional random walk
whose expected length depends on how densely the graph is connected and how
good the decision procedure is. Each regime of the walk was connected to a
phenomenon: drift toward the goal to improvement with practice, a balance
of branching against the shortening of paths to the fan effect. In this
sense the 1988 paper already had the explanatory shape I have since been
arguing for in CogLean \citep{coglean,cogleanintro}: mechanisms, together
with stated conditions, entail phenomena, including conditions under which
a phenomenon fails or reverses.

What it did not have was a way to check that its conclusions followed. The
derivations were done by hand, the approximations were judged by eye, and
the agreement between theory and simulation was read off log-log plots. So
I have written the 1988 theory and the bit game in Lean
\citep{demoura2021,mathlib2020}, where the definitions that run are the
same objects that theorems are about, and every claim of the paper has been
given a status: proved, sharpened, reinterpreted, or refuted
(Table~\ref{tab:claims}). The work is in a public repository
\citep{shh88revisited}; every result below names the theorem that
establishes it.

The outcome is a mixed verdict, which is the interesting kind. The
mathematics of the walk is right. But the two conclusions the paper drew
from it about practice do not follow: one is contradicted by the paper's
own idealized model, and the other names the wrong curve. And the place
where theory and game come apart can be located exactly.

\section{The model}
\label{sec:model}

I restate the 1988 model briefly, as the formalization reads it; the full
list of postulates, conditions and idealizations, with page references, is
in \code{notes/POSTULATES.md}.

\paragraph{The bit game.} A state is a vector of $B$ bits. An operator is a
pattern--result pair with don't-care positions: \code{?1?1?} $\to$
\code{?0?0?} applies to any state with 1s in the second and fourth
positions and sets them to 0. Play begins with the $2B$ one-bit operators,
which connect the states as a hypercube. A trial draws a start and a goal,
and performance is the number of operators applied to reach the goal
(\lean{SHH.BitGame.apply}, \lean{SHH.BitGame.initialOps}).

\paragraph{Learning.} After each trial, the operator that solves that game
in one step is added, with a don't-care wherever start and goal agree.
For \code{01110} $\to$ \code{10101} it is \code{01?10} $\to$ \code{10?01}
(\lean{SHH.BitGame.chunk\_paper}), which ``connects two pairs of states'':
it applies to exactly two of the 32 states
(\lean{SHH.BitGame.chunk\_paper\_applies}). In general a learned operator
takes its own start to its own goal (\lean{SHH.BitGame.chunk\_apply}).

\paragraph{The decision procedure.} At a state, the applicable operators
whose results are closest to the goal (by Hamming distance) are good, the
rest bad; each bad one is discarded with probability $r$, and one operator
is chosen at random from what remains. So $r = 1$ is a perfect procedure
and $r = 0$ a random walk. The formalization computes the exact
distribution over next states, with the probabilities of good and bad
choices as binomial sums (\lean{SHH.BitGame.stepDist}).

\paragraph{The theory.} To analyze typical rather than particular problem
spaces, the paper idealizes the graph as random, with $n$ nodes and mean
branching $\mu$, and then as a tree: start and goal are about the diameter
$D = \ln n/\ln\mu$ apart, and at each node one edge leads toward the goal
and $\mu - 1$ away. The decision procedure leaves $(\mu-1)(1-r)$ effective
bad choices, so a step goes toward the goal with probability
\begin{equation}
p_- = \frac{1}{1 + (\mu - 1)(1 - r)} .
\label{eq:pminus}
\end{equation}
The search becomes a walk on distance $0,\dots,D$ from the goal
(\lean{SHH.GD.pOf}, \lean{SHH.GD.DOf}). Everything in the theory depends
on this idealization, and the formalization keeps it separate from the
game, as an explicit postulate. That separation turns out to matter.

\section{The theory, checked}
\label{sec:theory}

\paragraph{The expected number of steps.} The paper gives the expected
number of steps $s$ to reach the goal from distance $D$ as its Eq.~2,
obtained by ``standard techniques'' \citep{vankampen1981}. The
formalization states the walk as a linear system for the expected steps
from every distance (\lean{SHH.Walk.Sys}), proves a closed form satisfies
it (\lean{SHH.Walk.term\_sys}), and proves the system has only one
solution (\lean{SHH.Walk.sys\_unique}); together these make the closed form
the expected number of steps (\lean{SHH.Walk.sys\_eq\_s}). The paper's Eq.~2
as printed,
\[
s = \frac{2}{(2p_- - 1)^2}\left[\,p_-^2 + p_-(D-1) - \frac{D}{2} + \frac{(1-p_-)^{D+1}}{p_-^{D-1}}\right],
\]
is exactly this function for every $p_-$ other than $0$ and $\tfrac12$
(\lean{SHH.Walk.eq2\_printed}). The 1988 derivation was right.

One detail had to be decided before any of this could be stated: what the
walk does at distance $D$, where the paper says only that it is
``constrained to remain within distance $D$''. The paper's own value for the
symmetric case, $s = D^2$, settles it: that value holds if the walk is
reflected back at $D$ (\lean{SHH.Walk.s\_half}), and not if it may stay
there, which gives $D^2 + D$ (\lean{SHH.Walk.stay\_half\_sys}). A choice the
paper left implicit is recoverable from what it reported.

\paragraph{The regimes.} When steps toward the goal are more likely than
steps away ($p_- > \tfrac12$), the paper's Eq.~3 says the walk drifts to the
goal in about $D/(p_- - p_+)$ steps. This holds as a two-sided bound: the
expected steps lie within $2p_-(1-p_-)/(2p_- - 1)^2$ below $D/(2p_- - 1)$,
a constant that does not grow with $D$ (\lean{SHH.Walk.drift\_bounds}). When
steps away are more likely, the paper's Eq.~4 gives the exponential term of
Eq.~2, and the rest is a term only linear in $D$
(\lean{SHH.Walk.eq4\_remainder}, \lean{SHH.Walk.s\_diffusion\_form}). In
every regime, more distance costs strictly more steps
(\lean{SHH.Walk.s\_strictMono\_D}).

\section{What the theory entails, and what it does not}
\label{sec:entails}

The paper then turned the walk into claims about practice. Learning adds
edges, so $\mu$ grows; more edges shorten the diameter, which helps, and add
bad choices, which hurts. Two of the paper's conclusions about which effect
wins do not follow.

\paragraph{A fan effect where the paper says there is none.} For a weak
decision procedure the paper approximates Eq.~4 by its Eq.~5, $s \approx
(2n/\mu)(1-r)^{\ln n/\ln\mu - 1}$, and concludes that ``when $r = 0$
(choices made at random), increasing the number of edges reduces the time
to solve the problem''. The approximation replaces $(\mu-1)^{D-1}$ by
$\mu^{D-1} = n/\mu$, dropping a factor $((\mu-1)/\mu)^{D-1}$ that is at most
$e^{-(D-1)/\mu}$ (\lean{SHH.GD.eq5\_error\_bound}). Where $\mu$ is small,
which is where learning starts, that factor is far from one. Restoring it
reverses the conclusion. With $n = 2^{16}$ nodes, where the diameter is an
integer for $\mu = 2, 4, 16, 256$ (\lean{SHH.GD.DOf\_pow2}), the paper's
own idealized model gives the expected steps in Table~\ref{tab:fan}: at
$r = 0$, raising the mean branching from 2 to 4 raises the expected steps
from 256 to 9{,}824, and only from $\mu = 16$ on do more edges help
(\lean{SHH.GD.fan\_pow2\_zero}). In the theory, there is a fan effect with
no decision procedure at all.

\begin{table}[H]
\centering\small
\begin{tabular}{@{}rrrrr@{}}
\toprule
$\mu$ & $D$ & $r = 0$ & $r = \tfrac12$ & $r = 1$ \\
\midrule
2   & 16 & 256     & 44.0    & 16 \\
4   & 8  & 9{,}824 & 255.5   & 8 \\
16  & 4  & 7{,}744 & 1{,}117.75 & 4 \\
256 & 2  & 512     & 257$^\dagger$ & 2 \\
\bottomrule
\end{tabular}
\caption{Expected steps in the paper's idealized model, $n = 2^{16}$, from
the proved closed form (\lean{SHH.GD.fan\_pow2\_zero},
\lean{SHH.GD.fan\_pow2\_half}, \lean{SHH.Walk.s\_one}). $^\dagger$From the
prototype (\code{tools/walk\_check.py}), not certified.}
\label{tab:fan}
\end{table}

\paragraph{A logarithm where the paper says power law.} With a strong
decision procedure the paper's Eq.~6 gives $s \approx \ln n/\ln\mu$ and
concludes that ``one obtains a power-law decay''. At $r = 1$ the
approximation is exact: the walk never steps away, so $s = D$
(\lean{SHH.Walk.s\_one}). To read this as a learning curve one needs a
further condition the paper leaves implicit, that each trial adds a bounded
number of edges, so that $\mu$ grows about linearly with trials. The curve
is then $T(t) = \ln n/\ln(\mu_0 + ct)$ (\lean{SHH.GD.curveGD}). It improves
with every trial (\lean{SHH.GD.curveGD\_improves}), and its proportional
gains shrink trial by trial (\lean{SHH.GD.curveGD\_diminishingRelative}),
the signature that distinguishes a power curve from an exponential one
\citep{heathcote2000}. Over trials 1--100 it lies within about 12\% of a
power law with exponent $-0.29$ (\code{tools/walk\_check.py}). But it is not
a power law. Its doubling ratio $T(2t)/T(t)$ tends to 1, so for any
exponent and tolerance, approximate scale invariance fails beyond some
trial (\lean{SHH.GD.curveGD\_not\_powerLaw}). The curve is logarithmic. On a
log-log plot over the one decade of Figure~1 of the 1988 paper, the
difference is invisible.

This does not undo the 1988 argument against building the power law into
the mechanism. It sharpens it. The graph-dynamic model does not build in a
power law, and, at least in its perfect-decision regime, it does not
produce one either. What it produces is a curve that cannot be told apart
from one over the range anyone plotted.

\section{The game itself}
\label{sec:game}

The theory's conclusions are about an idealization. Is the fan effect at
$r = 0$ a property of search in the bit game, or of the idealization? The
formalization can answer this exactly for small games, because the expected
number of steps in the game is the solution of a finite linear system, and
that solution is unique whenever the goal is reachable
(\lean{SHH.Markov.hit\_unique}, proved by the maximum principle). So a table
of values checked to satisfy the system is the expected number of steps.

On the three-bit game, adding one learned operator lowers the mean expected
steps over all 56 games at $r = 0$ and at $r = 1$, and raises it at
$r = \tfrac12$ (Table~\ref{tab:game}; \lean{SHH.BitGame.figure1\_pattern\_three}).
The prototype finds the same for every one of the 32 operators that could
be learned. This is the pattern of the 1988 simulations: improvement with
no decision procedure and with a perfect one, degradation with a mediocre
one. So the bit game does not show the fan effect at $r = 0$ that its
idealization does. The discrepancy is in the idealization: a hypercube is
full of short cycles, and the tree that the theory substitutes for it has
none. (Individual games can still slow down at $r = 0$: \code{000}
$\to$ \code{001} goes from 7 to $133/15$ expected steps,
\lean{SHH.BitGame.chunk\_slows\_game}. It is the average that improves.)

\begin{table}[H]
\centering\small
\begin{tabular}{@{}lccc@{}}
\toprule
& $r = 0$ & $r = \tfrac12$ & $r = 1$ \\
\midrule
initial operators & $58/7 \approx 8.286$ & $128/35 \approx 3.657$ & $12/7 \approx 1.714$ \\
plus the chunk \code{00?} $\to$ \code{11?} & $2603/315 \approx 8.263$ & $\approx 3.682$ & $23/14 \approx 1.643$ \\
\bottomrule
\end{tabular}
\caption{Mean expected steps over the 56 games of the three-bit game,
exact (\lean{SHH.BitGame.mean\_r0\_Base} and its five companions).}
\label{tab:game}
\end{table}

Two results about the perfect decision procedure complete the picture. With
the initial operators, every move it can make brings the state one bit
closer to the goal, for any number of bits, so it never searches or
backtracks (\lean{SHH.BitGame.r1\_moves\_closer}). Once operators have been
learned, this fails: on the four-bit game, with the learned operator
\code{0000} $\to$ \code{1111}, the game \code{0001} $\to$ \code{1110} has a
three-move solution (\lean{SHH.BitGame.three\_step\_path}) but the perfect
procedure expects $15/4$ moves (\lean{SHH.BitGame.greedy\_not\_optimal}).
Choosing the move that looks best by distance can miss the shortcut that
learning created. The 1988 paper said as much, that its $r = 1$ procedure
``only approximates optimal problem solving''; here is a case, exactly.

\begin{table}[t]
\centering\small
\begin{tabularx}{\linewidth}{@{}>{\raggedright\arraybackslash}p{0.42\linewidth}>{\raggedright\arraybackslash}X@{}}
\toprule
Claim of the 1988 paper & Status \\
\midrule
Eq.~2, the expected steps & proved, exactly as printed (\lean{SHH.Walk.eq2\_printed}) \\
$s = D^2$ when $p_+ = p_-$ & proved; fixes reflection at $D$ (\lean{SHH.Walk.s\_half}) \\
Eq.~3, drift & proved as a two-sided bound (\lean{SHH.Walk.drift\_bounds}) \\
Eq.~4, diffusion & proved as Eq.~2's exponential term (\lean{SHH.Walk.eq4\_remainder}) \\
Eq.~5 and ``at $r = 0$ more edges reduce the time'' & refuted in the paper's own idealized model (\lean{SHH.GD.fan\_pow2\_zero}) \\
Fan effect at intermediate $r$ & proved, theory (\lean{SHH.GD.fan\_pow2\_half}) and game (\lean{SHH.BitGame.figure1\_pattern\_three}) \\
Eq.~6 and ``a power-law decay'' & reinterpreted: logarithmic, not a power law (\lean{SHH.GD.curveGD\_not\_powerLaw}) \\
Eq.~2 ``produces the range of behaviors observed'' & fails at $r = 0$ because of the tree idealization (\lean{SHH.BitGame.figure1\_pattern\_three}) \\
$r = 1$ approximates optimal search & proved optimal on the initial operators (\lean{SHH.BitGame.r1\_moves\_closer}); not optimal with learned ones (\lean{SHH.BitGame.greedy\_not\_optimal}) \\
\bottomrule
\end{tabularx}
\caption{The 1988 claims and their status. The full list is in \code{notes/CLAIMS.md}.}
\label{tab:claims}
\end{table}

\section{Discussion}
\label{sec:disc}

The 1988 paper was mathematically careful, and most of its mathematics
survives a machine check unchanged. What did not survive were two inferences
from the mathematics to the psychology: an approximation applied where it
is poor, and a curve called by the name of the curve it resembles. Neither
is an unusual kind of mistake. Both are the kind a proof assistant catches
almost incidentally, because it will not accept ``approximately'' without
saying approximately what, over which range, and why.

The more interesting finding is about idealization. The 1988 paper moved
from the bit game to a random graph to a tree to a walk in a few sentences,
and treated the walk's behavior as the game's. In the formalization the
game and the walk are separate objects, and the step between them is a
postulate with a name. That is what made it possible to ask whether a
conclusion of the theory holds of the game, and to find one that does not.
I think this is the general lesson. Idealizations are where theories make
their least visible commitments; writing them down as first-class objects,
rather than as a sentence of motivation, makes them something one can test.

It also says something about the 1988 argument with ACT*. That argument was
that a theory should not get the power law of practice by building it into a
mechanism. The graph-dynamic model meets that standard, but the curve it
yields in its cleanest regime is logarithmic, and the evidence for the power
law has itself been contested \citep{heathcote2000}. The defensible claim
now is narrower and, I think, more interesting: that improvement with
diminishing relative returns, which is what the data most clearly show,
emerges from search dynamics without any learning curve being assumed, and
that the same dynamics, with a mediocre decision procedure, make practice
hurt.

\paragraph{Limitations.} The game results are for three and four bits, and
for one learning step at a time; the 1988 simulations used five bits and a
hundred trials of learning. Reproducing Figure~1 of that paper by
simulation, with the sampler checked against the exact small cases, is the
next step. The tree idealization is stated, not derived from random-graph
theory \citep{bollobas1985}, and the exact fan-effect results in the theory
use settings where the diameter is an integer. Nothing here concerns
people: the phenomena are properties of the model's behavior, and
connecting them to human practice data needs the usual assumptions about
what is measured.

\section{Conclusion}

The 1988 paper argued that learning need not improve performance, and that
whether it does depends on how decisions are made. Revisiting it with a
proof assistant confirms that argument and corrects two of its supporting
claims, and the corrections came from writing the theory down, not from new
mathematics. That is the point of the title. A model that only runs can
show that a behavior occurs; a model one can reason about shows what it
entails, where its approximations hold, and which of its idealizations is
responsible for a result. Thirty-eight years on, the graph-dynamic model
has more to say than its authors knew.

\section*{Availability}

All definitions and theorems are in SHH88Revisited \citep{shh88revisited}
(Lean 4.35.0, Mathlib v4.35.0). Every theorem is fully proved, with no
\code{sorry}, and depends only on Lean's standard axioms. The Python
prototypes in \code{tools/} have the same semantics as the Lean
definitions and generate the exact certificates.

\section*{Acknowledgements}

The original work was done with Tad Hogg and Bernardo Huberman at Xerox
PARC. Almost all of the Lean programming was accomplished by Claude Code
with extensive guidance from the author. Papers and presentations were
initially drafted by LLMs and then substantially rewritten by the author.
Any remaining errors, including hallucinations, are the author's alone.

\bibliographystyle{unsrtnat}
\bibliography{refs}

\end{document}
